\section{Problem Summary}
Given coin denominations and a target amount, return the minimum number of coins needed to make that amount. If it cannot be formed, return \(-1\).

Because each coin may be used repeatedly, this is not a one-time selection problem. It is a repeated-choice optimization problem over smaller amounts.

\section{Recognition Pattern}
This is a \textbf{1D dynamic programming} problem with an \textbf{unbounded knapsack-style transition}.

In interviews, recognize it when:
\begin{itemize}
\item the answer for amount \(x\) depends on answers for smaller amounts
\item choices can be reused multiple times
\item the question asks for a minimum or maximum count, not just feasibility
\end{itemize}

\section{Key Idea}
Let \(dp[x]\) be the minimum number of coins needed to form amount \(x\).

For every amount \(x\), try each coin \(c\):
\[
dp[x] = \min(dp[x], dp[x - c] + 1)
\]
whenever \(x - c \ge 0\) and \(dp[x - c]\) is reachable.

The recurrence is natural: if the last coin used is \(c\), then everything before that last step is exactly the subproblem for amount \(x - c\).

\section{Step-by-step Reasoning}
The brute-force recursive solution tries every coin choice at every step. That quickly becomes exponential because the same remaining amount is solved again and again.

Memoization fixes the repeated work, but the bottom-up version is usually easier to explain in an interview because the state is simple and the implementation is compact.

The important observation is that amounts from \(0\) up to the target form a valid build order. Once smaller amounts are solved, each larger amount can reuse those answers immediately.

Using a large sentinel value lets the code distinguish between reachable and unreachable states cleanly.

\section{Complexity}
\begin{itemize}
\item Time: \(O(amount \cdot n)\), where \(n\) is the number of coin types
\item Space: \(O(amount)\)
\end{itemize}

\section{Common Pitfalls}
\begin{itemize}
\item If you use a large sentinel such as `amount + 1`, do not add one to an unreachable state without checking it first.
\item Remember that `amount = 0` should return zero coins.
\item Mixing this problem with the counting version (\textbf{Coin Change II}) can lead to the wrong recurrence or the wrong interpretation of the DP state.
\end{itemize}

\section{Variants / Follow-ups}
This same template appears in \textbf{Coin Change II}, \textbf{Perfect Squares}, and several minimum-step or minimum-cost problems over integer states.

The main reusable lesson is to define a DP state that represents a smaller prefix of the answer space and then ask what the final move could have been.
